% titel: Traegheitsmoment von Molekuelen
% skript: Kap. Rotation, Traegheitsmoment mehrerer Massenpunkte

Ein planares dreiatomiges Molekül besteht aus drei punktförmigen Massen $m$
an den Eckpunkten eines gleichseitigen Dreiecks (Kantenlänge $d$). Das
Trägheitsmoment um verschiedene Achsen wird einmal für punktförmige, einmal
für kontinuierlich (gleichmässig) auf der Dreiecksfläche verteilte Masse
bestimmt.

\begin{werkzeugkasten}[Werkzeuge]
\textbf{Trägheitsmoment diskreter Massenpunkte} (Abstand $r_i$ von der Achse):
\begin{equation}
  I = \sum_i m_i\, r_i^2.
\end{equation}

\textbf{Kontinuierliche (konstante) Massenverteilung} mit Flächendichte
$\sigma=m/A$:
\begin{equation}
  I = \sigma \int_A r_\perp^2\,\mathrm dA.
\end{equation}

\textbf{Gleichseitiges Dreieck:} Kantenlänge $d$, Höhe
\begin{equation}
  h = \frac{\sqrt3}{2}\,d,
  \qquad
  A = \frac{\sqrt3}{4}\,d^2,
\end{equation}
Schwerpunkt $S$ auf der Winkelhalbierenden, Abstand $\frac13 h$ von einer
Seite; die Winkelhalbierende durch eine Ecke teilt die gegenüberliegende Seite
im rechten Winkel (die beiden anderen Massen haben von ihr den lotrechten
Abstand $d/2$).

\textbf{Achsen:} $A1$ liegt \emph{in einer Seitenlinie} (durch zwei Atome),
$A2$ ist die \emph{Winkelhalbierende durch den Schwerpunkt} (senkrecht zur
gegenüberliegenden Seite, durch ein Atom).
\end{werkzeugkasten}

\begin{werkzeugkasten}[Vorgehen in allgemeiner Form]
Punktförmig: Abstände der drei Massen von der jeweiligen Achse ablesen und
$I=\sum m r^2$. Kontinuierlich: $\sigma=3m/A$ und das Flächenintegral von
$r_\perp^2$ über das Dreieck für jede Achse (Achse in der Seitlinie bzw.
Winkelhalbierende) auswerten.
\end{werkzeugkasten}

\begin{achtung}[Typische Fallen]
\begin{itemize}
  \item Für die Achse $A1$ wie für $A2$ liegen \emph{je nach Achse andere}
       Atome auf bzw.\ in der Achse: $r=0$ für Atome in der Achse.
  \item Bei der kontinuierlichen Verteilung bleibt die \emph{Gesamtmasse}
       $3m$ erhalten: $\sigma=3m/A$, nicht $m/A$.
  \item Die Höhe des gleichseitigen Dreiecks ist $\frac{\sqrt3}{2}d$, nicht
       $d$.
\end{itemize}
\end{achtung}
